\documentclass[10pt]{article}

\usepackage[utf8]{inputenc}

\usepackage[T1]{fontenc}

\usepackage{amsmath}

\usepackage{amsfonts}

\usepackage{amssymb}

\usepackage[version=4]{mhchem}

\usepackage{stmaryrd}

\usepackage{graphicx}

\usepackage[export]{adjustbox}

\graphicspath{ {./images/} }



\begin{document}

можно считать составленными из чисел 0 и 1 (обычно под «тета-характеристиками» подразумевают именно такие характеристики). Для Т.-ф. с тета-характеристикой $H$ формулы (5) принимают вид



\[

\begin{gathered}

\theta_{H}\left(z+e_{\mu}\right)=(-1)^{h_{\mu}} \cdot \theta_{H}(z), \\

\theta_{H}\left(z+e_{\mu} A\right)=(-1)^{h_{\mu}^{\prime}} \exp \left[-\pi i\left(a_{\mu \mu}+2 z_{\mu}\right)\right] \cdot \theta_{H}(z) .

\end{gathered}

\]



Тета-характеристика $H$ наз. четной или нечетной в соответствии с тем, четна или нечетна Т.-ф. $\theta_{H}(z)$. Другими словами, тета-характеристика $H$ четная или нечетная в зависимости от того, четное или нечетное число $h^{\prime} h^{\top}$, поскольку



\[

\theta_{H}(-z)=(-1)^{h^{\prime} h^{\top}} \cdot \theta_{H}(z)

\]



Всего различных тета-характеристик $2^{2 p}$, из них $2^{p-1}\left(2^{p}+1\right)$ четных и $2^{p-1}\left(2^{p}-1\right)$ нечетных. Т.-ф. $\theta_{H}(z)$ обращается в нуль в тех точках $\left(g^{\prime}+g A\right) / 2$, тетахарактеристики к-рых $G=\left\|_{g}^{g},\right\|$ в сумме с $H$ составляют нечетную тета-характеристику. К. Якоби в своей теории әллиптич. функций использовал именно Т.-ф. с полуцелыми характеристиками, только периоды у него равнялись $\pi i$, а не 1.



Пусть $k$ - натуральное число. Целая трансцендентная функция $\theta_{\Gamma}(z)$ наз. Т.-ф. п о р я д к а $k$ с $\mathbf{x}$ a-1 р а к т е р и с т к о й $\Gamma$, если для нее выполняются тождества



\[

\begin{gathered}

\theta_{\Gamma}\left(z+e_{\mu}\right)=\exp \left(2 \pi i \gamma_{\mu}\right) \cdot \theta_{\Gamma}(z), \\

\theta_{\Gamma}\left(z+e_{\mu} A\right)=\exp \left[-\pi i\left(k a_{\mu \mu}+2 k z_{\mu}-2 \gamma_{\mu}^{\prime}\right)\right] \cdot \theta_{\Gamma}(z) .

\end{gathered}

\]



Напр., произведение $k$ Т.-ф. 1-го порядка есть Т.-ф. порядка $k$.



При помощи этих Т.-ф. 1-го порядка с полуцелыми характеристиками строятся мероморфные абелевы функции с $2 p$ периодами. Периоды произвольной абелевой функции от $p$ комплексных переменных удовлетворяют соотношениям Римана - Фробениуса, к-рые обеспечивают сходимость рядов, определяющих Т.-ф. с соответствующей системой модулей. А согласно теореме, сформулированной К. Вейерштрассом (К. Weierstrass) и доказанной А. Пуанкаре (Н. Poincaré), абелева функция может быть представлена в виде отношения целых Т.-ф. с соответствующей системой модулей. Для решения Я коби проблемы обращения абелевых интегралов строятся специальные Римана тета-бункции, аргументом к-рых является система точек $w_{1}, \ldots, w_{p}$ на римановой поверхности.



См. также Тета-ряд.



Лum.: [1] ч е б о т а p е в Н. Г., Теория алгебраических



\includegraphics[max width=\textwidth, center]{2023_07_09_52eacd0cbdccb74c34ccg-1}

z e r A., Lehrbuch der Theta-funktionen, Lpz., 1903; [4] C o nf o r t o F., Abelsche Funktionen und algebraische Geometrie,



ТЕТРАЦИКЛИЧЕСКИЕ КООРДИНАТЬІ Толомец. н а п ло с к о с т и- четыре числа $x_{1}, x_{2}, x_{3}, x_{4}$, подчиненные равенствам $\sigma x_{i}=k_{i} S_{i}, i=1,2,3,4$, где $S_{i}-$ степень точки относительно данных четырех окружностей, $k_{i}-$ произвольно заданные постоянные, $\sigma-$ множитель пропорциональности. Т. к. связаны соотношением 2-й степени, к-рое приводится к виду $x_{1}^{2}+x_{2}^{2}+$ $+x_{3}^{2}-x_{4}^{2}=0$, если исходные окружности взять ортогональными (из них три обязательно имеют действительные радиусы $\rho_{i}, i=1,2,3$, и одна - мнимый $\rho_{4}$ ), а числа $k_{i}$ равными $1 / \rho_{i}, i=1,2,3$, и $k_{4}=1 / i \rho_{4}$. Если в плоскости ввести декартовы координаты $\xi, \eta$, а в качестве трех действительных кругов взять $\xi=0, \eta=0$ (круги, проходящие через бесконечно удаленную точку плоскости), круг $\xi^{2}+\eta^{2}=1$ и мнимый круг $\xi^{2}+\eta^{2}=-1$, то тогда Т. к. точки на плоскости выразятся через декартовы координаты следующим образом:



\[

\sigma\left(x_{1}\right)=\xi, \quad \sigma x_{2}=\eta, \quad \sigma x_{3}=\frac{1-\left(\xi^{2}+\eta^{2}\right)}{2}, \quad \sigma x_{4}=\frac{1+\left(\xi^{2}+\eta^{2}\right)}{2} .

\]



Можно ввести Т. к. и для круга на плоскости. При указанном специальном выборе четырех основных кругов круг с центром в точке $\left(\xi_{0}, \eta_{0}\right)$ и радиусом $R_{0}$ имеет Т. к. $y_{i}, i=1,2,3,4$, определенные формулами



\[

\begin{gathered}

\sigma y_{1}=\xi_{0}, \sigma y_{2}=\eta_{0}, \quad \sigma y_{3}=\frac{1-\left(\xi_{0}^{2}+\eta_{0}^{2}-R_{0}^{2}\right)}{2}, \\

\sigma y_{4}=\frac{1+\left(\xi_{0}^{2}+\eta_{0}^{2}-R_{0}^{2}\right)}{2} .

\end{gathered}

\]



Т. к. точек и кругов на плоскости можно ввести с помощью стереографической проекции. При этом Т.к. точки на плоскости - однородные координаты соответствующей при стереографич. проектировании точки на сфере. Т. к. круга на плоскости - однородные координаты точки пространства, являющейся полюсом плоскости круга на сфере, соответствующего в стереографич. проекции кругу на плоскости, относительно этой сферы.



Обобщением Т. к. на случай трехмерного пространства являются пентасферические координаты.



Лит.: [1] К л е й н Ф., Высшая геометрия, пер. с нем., M.-Л., 1939; [2] Б у шा м а н о в а Г. В., Н о р д е н А. П., Элементы конформной геометрии, Казань, 1972. .



ТЕТРАӘДР - один из пяти типов правильных многогранников. Т. имеет 4 грани (треугольные), 6 ребер, 4 вершины (в каждой вершине сходится 3 ребра). Если $a$ - длина ребра Т., то его объем



\[

v=a^{3} \sqrt{2} / 12 \approx 0,1179 a^{3} .

\]



Т. является правильной треугольной пирамидой. БСЭ-з.



\begin{center}

\includegraphics[max width=\textwidth]{2023_07_09_52eacd0cbdccb74c34ccg-1(1)}

\end{center}



ТЕТРАӘДРА ПРОСТРАНСТВО - трехмерное пространство, являющееся пространством орбит действия бинарной групшы тетраэдра на трехмерной сфере. Эта группа определяется образующими $R, S$ и соотношениями $R^{2}=S^{3}=(R S)^{3}=E$. $M$. И. Войцеховский.



ТЕТРАЭДРАЛЫНЫЕ КООРДИНАТЫ Т о ч $\kappa$ и $P$ в т р е х м е р о м п р ос т р а н с т в е - числа $x_{1}$, $x_{2}, x_{3}, x_{4}$, пропорциональные (с заданным коэфффициентом пропорциональности) расстояниям от точки $\boldsymbol{P}$ до граней фиксированного тетраәдра. Аналогично вводятся общие Т. к. для любой размерности. Двумерный аналог Т. к. наз. т р е у голь пы м и к о о р ди нат а м и. См. также Баричентрические координаты.



ТИПИЧНО ВЕЩЕСТВЕННАЯ ФУНКЦИЯ̆ В о бл а с т и $B$ - функция $f(z)$, аналитическая в нек-рой области $B$ плоскости $z$, содержащей отрезки вещественной оси, если она вещественна на этих отрезках и $\operatorname{Im} f(z) \cdot \operatorname{Im} z>0$ при $\operatorname{Im} z \neq 0$. Основной класс Т. в. ф.к л а с с $T$ функций



\[

f(z)=z+\sum_{n=2}^{\infty} c_{n} z^{n}

\]



регулярных и типично вещественных в круге $|z|<1$ (cм. [1]). Из определения класса $T$ следует, что $c_{n}, n \geqslant 2$, вещественны. Класс $T$ содержит к л а с с $S_{r}$ функций



\[

f(z)=z+\sum_{n=2}^{\infty} c_{n} z^{n}

\]



с вещественными коэффициентами $c_{n}$, регулярных и однолистных в $|z|<1$. Если $f \in T$, то



\[

\varphi(z)=\frac{1-z^{2}}{z} f(z) \in C_{r},

\]



и, обратно, если $\varphi \in C_{r}$, то



\[

f(z)=\frac{z}{1-z^{2}} \varphi(z) \in T,

\]





\end{document}